% GENERATED FILE. Edit canonical lessons, then run npm run generate:books.
% canonical-source-hash: fnv1a64:10ebe9724f711025
% canonical-lessons: FR-C434-orienter, FR-C434-boussole, FR-C434-reperer, FR-C434-repere, FR-C434-egarer

\chapter{Trouver son chemin}
\label{ch:fr-trouver-son-chemin}
\hlchaptermodality{434}

\begin{chapteropening}
\textbf{By the end of this chapter:} \emph{I can find my bearings: getting oriented, the compass, spotting a landmark, the landmark itself and getting lost.}

Complete the last lesson of chapter 434: I can find my bearings: getting oriented, the compass, spotting a landmark, the landmark itself and getting lost.
\end{chapteropening}

\section[orienter]{orienter — to get one’s bearings (mostly reflexive: s’orienter); to guide}

\label{lesson:FR-C434-orienter}



\begin{quote}
Before the new one: say the French for the compensation, then the French for to pay damages to.
\end{quote}

\subsection*{You'll want to know: orienter}
\textbf{orienter} — \textquotedblleft{}to get one’s bearings (mostly reflexive: s’orienter); to guide\textquotedblright{}.

\begin{quote}
Dans la vieille ville, je m'oriente mal.
\end{quote}

\begin{quote}
\emph{In the old town I find it hard to get my bearings.}
\end{quote}

Mostly reflexive: \emph{s'orienter}. \emph{Le sens de l'orientation} — a sense of direction.

\begin{cousinweb}
\emph{orienter} is built on \emph{orient} 'east', from Latin \emph{oriēns} 'rising', said of the sun, from \emph{orīrī} 'to rise' (Lewis \& Short): first to turn towards the east.
\end{cousinweb}

\begin{grammarlens}[title={What you've built}]
One more word for telling, explaining and describing at B1.
\end{grammarlens}

\subsection*{Your turn}
\begin{itemize}
\raggedright
  \item \emph{Say it:} \emph{orienter}
  \item \emph{Say it:} \emph{orienter}, once more
  \item \emph{Say it:} a sentence of your own with \emph{orienter}
  \item \emph{Recall:} say the French for the compensation, then the French for to pay damages to, then say \emph{orienter} again
\end{itemize}

\begin{tcolorbox}[breakable,colback=teal!4,colframe=teal!35!black,title={Before you move on}]
What does orienter mean? (\textquotedblleft{}To get one’s bearings (mostly reflexive: s’orienter); to guide\textquotedblright{}.) Say it once more.
\end{tcolorbox}
\section[la boussole]{la boussole — the compass}

\label{lesson:FR-C434-boussole}



\begin{quote}
Before the new one: say the French for to pay damages to, then the French for to get one’s bearings.
\end{quote}

\subsection*{You'll want to know: la boussole}
\textbf{la boussole} — \textquotedblleft{}the compass\textquotedblright{}.

\begin{quote}
En forêt, il garde toujours sa boussole sur lui.
\end{quote}

\begin{quote}
\emph{In the woods he keeps his compass on him the whole time.}
\end{quote}

\emph{Perdre la boussole} — to lose one's head.

\begin{cousinweb}
\emph{boussole} was borrowed from Italian \emph{bussola} 'compass' (EtymDB), from Late Latin \emph{buxula} 'small box', built on \emph{buxus} 'box tree'.
\end{cousinweb}

\begin{grammarlens}[title={What you've built}]
One more word for telling, explaining and describing at B1.
\end{grammarlens}

\subsection*{Your turn}
\begin{itemize}
\raggedright
  \item \emph{Say it:} \emph{la boussole}
  \item \emph{Say it:} \emph{la boussole}, once more
  \item \emph{Say it:} a sentence of your own with \emph{la boussole}
  \item \emph{Recall:} say the French for to pay damages to, then the French for to get one’s bearings, then say \emph{la boussole} again
\end{itemize}

\begin{tcolorbox}[breakable,colback=teal!4,colframe=teal!35!black,title={Before you move on}]
What does la boussole mean? (\textquotedblleft{}The compass\textquotedblright{}.) Say it once more.
\end{tcolorbox}
\section[repérer]{repérer — to spot; to find one’s way (se repérer)}

\label{lesson:FR-C434-reperer}



\begin{quote}
Before the new one: say the French for to get one’s bearings, then the French for the compass.
\end{quote}

\subsection*{You'll want to know: repérer}
\textbf{repérer} — \textquotedblleft{}to spot; to find one’s way (se repérer)\textquotedblright{}.

\begin{quote}
Elle a repéré un petit café près de la gare.
\end{quote}

\begin{quote}
\emph{She spotted a little café near the station.}
\end{quote}

\emph{Se repérer} — to find your way around.

\begin{cousinweb}
\emph{repérer} is built on \emph{repère} 'mark, landmark' (EtymDB).
\end{cousinweb}

\begin{grammarlens}[title={What you've built}]
One more word for telling, explaining and describing at B1.
\end{grammarlens}

\subsection*{Your turn}
\begin{itemize}
\raggedright
  \item \emph{Say it:} \emph{repérer}
  \item \emph{Say it:} \emph{repérer}, once more
  \item \emph{Say it:} a sentence of your own with \emph{repérer}
  \item \emph{Recall:} say the French for to get one’s bearings, then the French for the compass, then say \emph{repérer} again
\end{itemize}

\begin{tcolorbox}[breakable,colback=teal!4,colframe=teal!35!black,title={Before you move on}]
What does repérer mean? (\textquotedblleft{}To spot; to find one’s way (se repérer)\textquotedblright{}.) Say it once more.
\end{tcolorbox}
\section[le repère]{le repère — the landmark, the reference point}

\label{lesson:FR-C434-repere}



\begin{quote}
Before the new one: say the French for the compass, then the French for to spot.
\end{quote}

\subsection*{You'll want to know: le repère}
\textbf{le repère} — \textquotedblleft{}the landmark, the reference point\textquotedblright{}.

\begin{quote}
L'église est un bon repère pour retrouver l'hôtel.
\end{quote}

\begin{quote}
\emph{The church is a good landmark for finding the hotel again.}
\end{quote}

\emph{Un point de repère} — a landmark.

\begin{cousinweb}
\emph{repère} is usually explained as a spelling of \emph{repaire} 'den, resting place' reshaped after Latin \emph{reperīre} 'to find' (Lewis \& Short); the details are not settled.
\end{cousinweb}

\begin{grammarlens}[title={What you've built}]
One more word for telling, explaining and describing at B1.
\end{grammarlens}

\subsection*{Your turn}
\begin{itemize}
\raggedright
  \item \emph{Say it:} \emph{le repère}
  \item \emph{Say it:} \emph{le repère}, once more
  \item \emph{Say it:} a sentence of your own with \emph{le repère}
  \item \emph{Recall:} say the French for the compass, then the French for to spot, then say \emph{le repère} again
\end{itemize}

\begin{tcolorbox}[breakable,colback=teal!4,colframe=teal!35!black,title={Before you move on}]
What does le repère mean? (\textquotedblleft{}The landmark, the reference point\textquotedblright{}.) Say it once more.
\end{tcolorbox}
\section[égarer]{égarer — to mislay; to get lost (s’égarer)}

\label{lesson:FR-C434-egarer}



\begin{quote}
Before the new one: say the French for to spot, then the French for the landmark.
\end{quote}

\subsection*{You'll want to know: égarer}
\textbf{égarer} — \textquotedblleft{}to mislay; to get lost (s’égarer)\textquotedblright{}.

\begin{quote}
Sans carte, ils se sont égarés dans les petites rues.
\end{quote}

\begin{quote}
\emph{Without a map they got lost in the little streets.}
\end{quote}

\emph{S'égarer} — to lose your way; \emph{égarer ses clés} — to mislay your keys.

\begin{cousinweb}
\emph{égarer} comes from Old French \emph{esgarer} (EtymDB), which is usually traced to a Frankish verb meaning 'to watch over', with \emph{es-} 'away'.
\end{cousinweb}

\begin{grammarlens}[title={What you've built}]
One more word for telling, explaining and describing at B1.
\end{grammarlens}

\subsection*{Your turn}
\begin{itemize}
\raggedright
  \item \emph{Say it:} \emph{égarer}
  \item \emph{Say it:} \emph{égarer}, once more
  \item \emph{Say it:} a sentence of your own with \emph{égarer}
  \item \emph{Recall:} say the French for to spot, then the French for the landmark, then say \emph{égarer} again
\end{itemize}

\begin{tcolorbox}[breakable,colback=teal!4,colframe=teal!35!black,title={Before you move on}]
What does égarer mean? (\textquotedblleft{}To mislay; to get lost (s’égarer)\textquotedblright{}.) Say it once more.
\end{tcolorbox}
